\subsection{The monoid lower bound: completed reductions and exact scope}
\label{A:monoid-lower-audit}

We distinguish a complete classification inside an explicitly specified
action hull from the classification of all finite monoids. This subsection
proves the three-row gates and the fixed-hull classification; the completed
large-band and four-by-four arguments follow in
Sections~\ref{A:large-band-target-proof} and \ref{A:monoid-global-answer}.
Together they establish the global answer and classification in
Theorem~\ref{A:monoid-minimum23-classification}.

\begin{lemma}[Proved: explicit coordinates on the minimal ideal]
\label{A:minimal-ideal-coordinates}
Let $K$ be the minimal ideal of a finite monoid $M$ and choose
an idempotent $e\in K$. Then $G=eMe$ is a group, and there
are finite sets $I,J$ identifying $K$ with $I\times G\times J$,
with multiplication
$(i,g,j)(i',g',j')=(i,gp_{ji'}g',j')$ for $p_{ji'}\in G$.
If $G=1$ this is a rectangular band. If $K=M$, then $M$ is a group.
\end{lemma}
\begin{proof}
The group assertion for $eMe$ is
Lemma~\ref{CP2:C-minimal-corner}; we give the coordinates explicitly.
Put $I=\{a\in Me:ea=e\}$ and
$J=\{b\in eM:be=e\}$. For $x\in Me$, the element $ex$
lies in $G$, and $x=a g$ uniquely with $g=ex$ and
$a=xg^{-1}\in I$. Dually $y\in eM$ is uniquely $h b$,
where $h=ye\in G$ and $b=h^{-1}y\in J$.
Minimality gives $K=MeM$, so every $z\in K$ can be written
$a g b$. Conversely, such a factorization recovers
$g=eze$, $a=ze\,g^{-1}$ and $b=g^{-1}ez$, proving uniqueness.
The product is
$a g b a'g'b'=a g(ba')g'b'$, and $ba'\in eMe=G$.
This proves the formula with $p_{ba'}=ba'$.
If $1\in MeM$, a factorization $1=men$ in a finite monoid
forces $e$ to be a unit: a one-sided inverse in a finite
monoid is two-sided, and a product is a unit only if each
factor is a unit. An idempotent unit is $1$, so $e=1$ and
$M=eMe=G$.
\end{proof}

\begin{lemma}[Proved: two augmentation shifts]
\label{A:band-Tor}
Let $M$ be a finite monoid with proper rectangular minimal ideal
$K=I\times J$, and let $k$ be $\Z$ or a field. Set
$R=kM$, $A=R/kK$, $V=\ker(kI\to k)$ as a left $R$-module,
and $W=\ker(kJ\to k)$ as a right $R$-module. There are natural
isomorphisms
\[
 H_{n+2}(BM;k)\cong\operatorname{Tor}^{A}_{n}(W,V)
 \quad(n\geq0).
\]
In particular $H_2(BM;k)$ is the quotient of $W\otimes_kV$ by
$wm\otimes v-w\otimes mv$ for $m\in M$.
\end{lemma}
\begin{proof}
Choose $e=(i_0,j_0)\in K$. Multiplication identifies the $R$-bimodule
$kK$ with $Re\otimes_k eR$: on the standard bases this is the
bijection $(i,j)\mapsto(i,j)$. Consequently $kK$ is projective on
both sides. Also $K$ annihilates $V$ and $W$, and
$kK\otimes_R V=Re\otimes_k eV=0$.
The right-projective resolution $0\to kK\to R\to A\to0$
therefore gives $\operatorname{Tor}^R_q(A,V)=0$ for $q>0$ and
$A\otimes_R V=V$. Filtering the two-sided double bar complex by
the $A$ resolution yields
\[
 E^2_{p,q}=\operatorname{Tor}^A_p
   (W,\operatorname{Tor}^R_q(A,V))
 \Longrightarrow\operatorname{Tor}^R_{p+q}(W,V).
\]
Only $q=0$ remains. Finally use the exact augmentation sequences
$0\to V\to Re\to k\to0$ and
$0\to W\to eR\to k\to0$.
Their middle terms are projective, and $eR\otimes_RV=eV=0$;
the long exact Tor sequences give the stated two shifts, including
$H_2=W\otimes_AV$. All constructions commute with homomorphisms
which preserve the identified band and its actions.
\end{proof}

\begin{lemma}[Proved: preservation of $H_2$ and surjection on $H_3$]
\label{A:band-quotient-H3}
Suppose $q:M\twoheadrightarrow N$ is a surjective monoid
homomorphism whose restriction identifies their rectangular minimal
bands bijectively. Then $q$ induces an isomorphism on $H_2(-;k)$
and a surjection on $H_3(-;k)$, for $k=\Z$ and for every field.
\end{lemma}
\begin{proof}
The identified row and column sets have the same actions: every
element of $N$ has a preimage, and multiplication on the band
determines the action. Their balanced presentations of $H_2$ are
therefore identical. For the second assertion use the natural
two-sided bar complexes computing the Tor groups in
Lemma~\ref{A:band-Tor}. The chain map is the identity in degree
zero, whose group is $W\otimes_kV$, and is surjective in degree
one because $kM/kK\to kN/kK$ is surjective. If $z$ is a target
one-cycle and $y$ is any chain lifting it, then $dy=0$ because
its image is $dz=0$ and the degree-zero map is injective. Thus
every homology class in degree one is represented by the image
of a source cycle. Naturality of the two shifts proves the claim.
This argument does not claim an isomorphism in higher degrees.
\end{proof}

\begin{lemma}[Proved: rational minimal-group collapse]
\label{A:minimal-group-collapse}
For a finite monoid whose minimal ideal has Rees coordinates
$I\times H\times J$, collapse the $H$ coordinate inside the
ideal and leave every other element distinct. The resulting
quotient $q:M\to\overline M$ is a monoid homomorphism, its
minimal ideal is $I\times J$, and $Bq$ is a rational homology
equivalence and hence a rational cohomology-ring isomorphism.
\end{lemma}
\begin{proof}
Left multiplication changes $(i,h,j)$ to
$(\sigma(i),a_i h,j)$ and right multiplication changes it to
$(i,hb_j,\tau(j))$. These descriptions follow by commuting with
all translations from the opposite side. They show directly
that forgetting $h$ is a congruence, including a nontrivial
Rees sandwich matrix.

Put $R=\mathbb QM$ and choose an idempotent $e$ in the minimal
ideal. The multiplication map identifies
$\mathbb QK$ with $Re\otimes_{\mathbb QH}eR$.
The module $Re$ is free over $\mathbb QH$ on the right, so
$\mathbb QK$ is projective as a right $R$-module. Averaging
the group coordinate in each $H$ class defines an $R$-bimodule
idempotent on $\mathbb QK$, since left and right translations
permute those coordinates. Its kernel is exactly
$J_0=\ker(R\to\mathbb Q\overline M)$, which is consequently
right projective. On $\mathbb QK\otimes_R\mathbb Q$ the
averaging operator is the identity; hence
$J_0\otimes_R\mathbb Q=0$. The resolution
$0\to J_0\to R\to\mathbb Q\overline M\to0$ now gives
vanishing higher Tor with the trivial module. The change of
rings argument used in Lemma~\ref{A:band-Tor} identifies the
homology through the specified map $q$. The induced map on
cohomology respects cups, proving the ring statement.
\end{proof}

\begin{lemma}[Proved: a rank-one balanced form has square zero]
\label{A:rank-one-square}
For a finite monoid with rectangular minimal band, every rational
degree-two class represented by a balanced matrix of rank at most
one has zero square. No finiteness assumption on $BM$ is needed.
\end{lemma}
\begin{proof}
Write the nonzero matrix as $L=b\otimes a$, with $a_0=b_0=0$.
Balance means $B_m^tL=LA_m$, where $A_m,B_m$ are the induced
augmentation operators. Equality of these two rank-one tensors
gives a common scalar character $\chi_m$ and affine laws
\[
 a_{\sigma_m(i)}=a_m+\chi_m a_i,\qquad
 b_{\tau_m(j)}=b_m+\chi_m b_j.
\]
Here $a_m=a_{\sigma_m(0)}$ and $b_m=b_{\tau_m(0)}$.
The composition rules are
$a_{mn}=a_m+\chi_m a_n$ and
$b_{mn}=b_n+\chi_n b_m$.
A nonzero slope permutes a finite nonconstant rational value set,
so it is $1$ or $-1$. A slope-one translation is zero, and all
slope-minus-one maps have a common center on each value set.
Choose those centers $A,B$, or zero if no reflection occurs, and set
\[
 a'_m=a_m+A(\chi_m-1),\qquad
 b'_m=b_m+B(\chi_m-1).
\]
Then $\chi_m a'_m=\chi_m b'_m=0$ for every $m$.
The affine two-cocycles $u(m,n)=b_m a_n$ and
$u'(m,n)=b'_m a'_n$ obey
\[
 u'-u=-\delta\bigl(Ab+Ba+AB(\chi-1)\bigr).
\]
For the normalized three-cochain
$w(m,n,p)=b'_m a'_n b'_n a'_p$, substitution in the five terms
of the coboundary gives
\[
 (\delta w)(m,n,p,q)
 =u'(m,n)u'(p,q)
   +b'_m a'_q\chi_n b'_n\chi_p a'_p
 =u'(m,n)u'(p,q).
\]
Thus $[u]^2=0$.

For completeness, balanced forms do give all rational degree-two
classes. By Lemma~\ref{A:band-Tor}, their vector space is the
dual of the balanced quotient of $W\otimes V$. The formula
$u_L(m,n)=L_{\tau_m(0),\sigma_n(0)}$, with zero row and
column at index zero, satisfies the cocycle equation by balance.
The band cycles
\[
 -[x_{ij}|x_{ij}]-[x_{00}|x_{00}]
 +[x_{i0}|x_{i0}]+[x_{ij}|x_{00}]
\]
pair with $u_L$ as $-L_{ji}$. These cycles are the images of
the elementary tensors under the two augmentation connecting
maps. The formula therefore realizes exactly the dual
identification, rather than just a subspace of cohomology.
\end{proof}

\begin{corollary}[Proved: a census-free exclusion through order nine]
\label{A:monoids-nine}
No finite monoid of order at most nine has classifying space
$\CP^2$. More generally, a minimal ideal with at most two rows
or at most two columns forces every rational degree-two square
to vanish, independently of the maximal group and of finite
homotopy type.
\end{corollary}
\begin{proof}
Apply Lemma~\ref{A:minimal-group-collapse}. With one row or
one column, one augmentation module is zero, so $H_2$ vanishes.
With two rows or two columns every balanced matrix has rank
at most one, and Lemma~\ref{A:rank-one-square} applies.
If both row and column counts are at least three, the minimal
ideal has at least nine elements. The identity lies outside a
proper minimal ideal, giving at least ten elements in the
monoid. If the minimal ideal contains the identity, the monoid
is a group; its positive rational homology vanishes. These
cases exhaust all monoids of order at most nine.
\end{proof}

\paragraph{Action hulls and the exact finite scope.}
For $r$ rows and $s$ columns, put $d=(r-1)(s-1)$. An actor is
a pair of transformations $(\sigma,\tau)$; its multiplication
is $(\sigma,\tau)(\alpha,\beta)=(\sigma\alpha,\beta\tau)$.
Let $R_{\sigma,\tau}\subset\mathbb Q^d$ be the space of
coefficient relations in $B_\tau^tL=LA_\sigma$.
For a proper sum $R$ of such spaces define
\[
 H(R)=\{(\sigma,\tau):R_{\sigma,\tau}\subset R\}.
\]
Balance is closed under multiplication, so this is a monoid.
If $R$ is generated by a monoid's actors, the hull contains all
of them and has precisely $R$ as its relation space. The action
homomorphism into $H(R)$ preserves the band and induces an
isomorphism on rational $H_2$, and a surjection on integral
$H_2$ whenever the band is unchanged. Thus a hypothetical
$\CP^2$ source has $\dim R=d-1$ and its degree-two generator
is pulled back rationally from that hull.

The new program \path{monoid/code/relation_lattice.py}
forms every actor relation space by rational row reduction and
closes proper sums modulo $S_r\times S_s$. Every new orbit is
marked in full. It suffices to join individual actor spaces to
each orbit representative, since the actor set is invariant
under the permutation group. A hyperplane has no larger proper
join, so those joins are settled by dimension, without a
numerical rank estimate. The independent checker
\path{monoid/code/relation_lattice_check.py} reconstructs the
relations as double differences of matrix entries and computes
permutation actions by expanding augmentation basis tensors.
It verifies every orbit, every proper join and every hull's
actor list and spanning equality. The resulting counts are
\[
\begin{array}{c|r|r|r|r}
 (r,s)&r^rs^s&\text{proper spaces}&S_r\times S_s\text{ orbits}
       &\dim H^2=1\text{ orbits}\\\hline
 (3,3)&729&111&13&6\\
 (3,4)&6912&2767&53&21\\
 (3,5)&84375&68719&214&71
\end{array}
\]
The $3\times3$ count is thirteen because opposition has not
been divided out; the inherited eleven-orbit count used it.

\begin{proposition}[Proved: the $3\times3$, $3\times4$, and
 $3\times5$ gates for $\CP^2$ through order twenty-three]
\label{A:three-row-classification}
No finite monoid of type $\CP^2$ has a minimal ideal with
three rows and three columns, with any maximal group. None has
a $3\times4$ rectangular minimal band. Among monoids of order
at most twenty-three with $3\times5$ rectangular minimal band,
the only $\CP^2$ model, up to isomorphism, is the monoid of
Theorem~\ref{A:monoid-upper23}. All three statements also hold
after opposition.
\end{proposition}
\begin{proof}
Only the one-dimensional balanced-form spaces in the preceding
table can occur. Their exact partition in our orbit numbering is
\[
\begin{array}{c|r|r|r|l}
 (r,s)&\operatorname{rank}L=1&\text{duplicate values}
 &\text{new hulls}&\text{new hulls with }H_4(\mathbb Q)\ne0\\\hline
 (3,3)&3&0&3&11\ (|H|=30)\\
 (3,4)&6&4&11&38\ (25),\quad40\ (23)\\
 (3,5)&8&37&26&164\ (34),\quad198\ (30),\quad200\ (30)
\end{array}
\]
The first column is covered by Lemma~\ref{A:rank-one-square}.
For duplicate values, identify column points on which every
balanced form agrees; balance preserves this equivalence
relation under all column actions. The action map descends to
the smaller set, and the displayed affine cocycles factor
through it. The augmentation quotient is surjective. Hence
the smaller action image has the same nonzero one-dimensional
rational balanced quotient and is covered inductively by the
smaller-shape argument. A repeated row is handled in the same
way, or by opposition. All source finiteness assumptions below
are imposed on the original monoid, not on this intermediate
image.

For each new hull the program
\path{monoid/code/hull_prefixes.py} constructs an integral
projective resolution prefix through degree five. Its projectives
are principal summands generated by explicitly checked idempotents
or differences of nested idempotents. The certificate stores their
bases, the entire action, differentials, and tensor differentials.
The separate checker
\path{monoid/code/projective_prefix_check.py} reconstructs the
table and every projective, checks all equivariance equations,
zero composites, rational ranks and integral Smith invariants.
Saturated images and the complementary kernel ranks prove exactness.
Thus its calculation of $H_4$ depends on an exact resolution,
not on a guessed truncation. Every new hull omitted from the
last column has verified zero rational fourth homology, which
kills a pulled-back square.

We explain the cyclic-character exclusions used for the few
remaining hulls. For an integral affine cocycle $u$ and an actor
$a$, let $G_a$ be its eventual cyclic group, $e$ its identity,
and $g=ea$. Its restricted degree-two class is
\[
 c_a=\sum_{h\in G_a}u(h,g)\pmod{|G_a|}.
\]
Indeed $u(e,h)=u(h,e)=u(e,e)$; subtracting this constant
normalizes the cocycle and changes the displayed sum by a
multiple of $|G_a|$. This is the elementary cyclic cocycle
formula for $H^2(BG_a;\Z)$.
If an actor with $c_a\ne0$ were the image of an element of a
finite monoid $M$, its eventual group would surject onto $G_a$.
The nonzero character would remain nonzero under inflation,
since $H^2(BG;\Z)=\operatorname{Hom}(G,\mathbb Q/\Z)$ for
finite groups and precomposition with a surjection is injective.
Such a restriction cannot come from a line bundle on a finite
CW model of $BM$. The reduced line-bundle class is nilpotent
by Lemma~\ref{CP2:C-Knilpotence}. Choose a cyclic prime-power
subgroup on which the restricted character is nontrivial.
Atiyah's completion theorem \cite[Theorem~7.2]{CP2-Atiyah1961}
identifies its K-theory with the completed representation ring.
Evaluation at a primitive prime-power root of unity maps the
reduced character to a nonzero element $\zeta^a-1$ of the
corresponding $p$-adic cyclotomic integer domain. The
augmentation ideal maps into its maximal ideal, so evaluation
extends to the completion, and no power of this nonzero
element vanishes. This contradicts nilpotence. A local group is interpreted as
the endomorphisms of its idempotent in the Karoubi envelope,
so the classifying-space map is well typed even when $e\ne1$.

The independent program
\path{monoid/code/character_prune_check.py} recomputes all
these residues, grows allowed submonoids by generated closures,
and compares every maximal target with its independently
checked integral projective prefix. The results are
\[
\begin{array}{c|r|r|l}
 \text{hull}&\text{forbidden actors}&\text{maximal allowed targets}
 &\text{verified conclusion}\\\hline
 (3,3),11&8&9\text{ of order }16&H_4(\mathbb Q)=0\\
 (3,4),38&2&1\text{ of order }23&H_4(\mathbb Q)=0\\
 (3,5),198&4&1\text{ of order }26&H_4(\mathbb Q)=0\\
 (3,5),200&4&1\text{ of order }26&H_4(\mathbb Q)=0
\end{array}
\]
For the first line the producer tested all $4096$ subsets,
finding $154$ closed sets; the independent closure-growth
checker recovers exactly these sets and the nine maxima.
In the last three lines the entire allowed set is already
closed, which proves directly that it is the unique maximum.
The degree-two class factors through one of these targets,
so its square vanishes on the original source.

The exceptional $3\times4$ hull 40 has order twenty-three.
The newly checked resolution prefix has abelian ranks
$(3,12,23,22,15,12)$ in degrees zero through five and
coinvariant ranks $(1,0,1,0,1,0)$.
The exact integral bar comparison in
\path{monoid/certificates/hull_3x4_40_cup.json} gives
\[
 H_2=\Z,\quad H_3=0,\quad H_4=\Z,\qquad x^2=2y,
\]
where $x$ and $y$ are primitive. The checker verifies the
comparison on projective generators and all their translates,
the cocycle on every triple, and evaluations on integral cycle
bases. These conclusions need only exactness through degree
four; we make no assertion here about an unverified periodic
tail. A band source with $H_2=\Z$ maps surjectively on integral
$H_2$ to this hull, so its primitive two-class is pulled back
up to sign. Its square is then even, which excludes $\CP^2$.
The same argument applies after identifying duplicate values,
because the corresponding augmentation map is still surjective.

It remains to treat the $3\times5$ hull 164. Its matrix, with
zero row and column suppressed, is
\[
 L=\begin{pmatrix}1&1\\0&1\\-1&0\\1&0\end{pmatrix}.
\]
It has thirty-four actors; four forbidden characters leave a
closed thirty-element allowed hull. The complete lattice of
closed subsets containing the identity and all fifteen band
elements has $2186$ members. Exactly $1711$ have order at most
twenty-three; among these $1160$ have integral $H_2=\Z$.
The two newly authored programs
\path{monoid/code/cp23_hull.py} and
\path{monoid/code/cp23_hull_check.py} obtain the same scope
by subset enumeration and by generated closures, respectively.
The latter reconstructs the integer balanced presentation,
its Smith form and the three two-sided-bar boundary matrices
over $\mathbb F_5$ for every relevant image. The resulting
counts are
\[
\begin{array}{c|r|r|r|r}
 |N|&\text{all images}&H_2(N;\Z)=\Z
  &\dim H_4(N;\mathbb F_5)=0
  &(b_2,b_3,b_4)=(1,0,1)\\\hline
16&1&0&0&0\\17&12&0&0&0\\18&60&0&0&0\\
19&166&12&12&0\\20&309&140&140&0\\
21&416&314&314&0\\22&421&380&376&0\\
23&326&314&310&2
\end{array}
\]
The four exceptions of order twenty-two have profile $(1,1,2)$;
the other two exceptions of order twenty-three have profile
$(1,2,3)$. They are excluded for every band-preserving source
by the surjection on $H_3(-;\mathbb F_5)$ in
Lemma~\ref{A:band-quotient-H3}. All images with zero fourth
homology over $\mathbb F_5$ have zero rational fourth homology,
so their pulled-back degree-two squares vanish.

The remaining two images are indices 1412 and 1651 in the
certificate. Both have the complete integral $\CP^2$ certificate
of Theorem~\ref{A:monoid-upper23}, and the displayed isomorphism
there identifies them. A source of order at most twenty-three
surjecting onto either twenty-three-element image must be
isomorphic to it. This proves the asserted classification.
For the $3\times3$ assertion with a nontrivial maximal group,
use Lemma~\ref{A:minimal-group-collapse}; the rational square
obstructions persist and all forbidden-character restrictions
are applied to the original finite-type source. This does not
assume that the collapsed quotient has finite CW type.
\end{proof}

\paragraph{Relation to the completed global proof.}
The preceding proposition classifies every monoid of order at most
twenty-three with a $3\times5$ rectangular minimal band, including
nonfaithful action sources. The $4\times4$ action computation is
completed independently in Proposition~\ref{A:four-by-four-through23}.
The remaining large bands are excluded by the target-specific
algebra argument of Section~\ref{A:large-band-target-proof}; it
avoids the stronger inherited arbitrary-module contracted-algebra
claim and its binary-pencil classification. The global minimum and
all monoids attaining it are therefore established in
Theorem~\ref{A:monoid-minimum23-classification}.
